% Extended Discussion Content to be added

\subsection{Mechanism Analysis}

\subsubsection{Why Does TRR Work?}

The effectiveness of TRR can be understood through multiple mechanisms:

\paragraph{Feature Co-activation as Style Signature}. Many audio effects produce characteristic patterns of co-activation across frequency bands. For example:
\begin{itemize}
	\item \textbf{Distortion}: Creates correlated harmonics that appear as structured off-diagonal patterns in the Gram matrix
	\item \textbf{Chorus}: Introduces periodic modulation that manifests as oscillatory patterns in time-averaged correlation structure
	\item \textbf{Compression}: Alters variance relationships, affecting diagonal and near-diagonal elements of $\mathbf{G}$
\end{itemize}

By capturing these second-order statistics, TRR provides a more discriminative signature of effect activation than first-order mean statistics alone.

\paragraph{Invariance to Irrelevant Variations}. TRR's aggregation across time and normalization provide robustness to:
\begin{itemize}
	\item \textbf{Recording level differences}: Normalization makes the representation invariant to overall signal energy
	\item \textbf{Temporal shifts}: Time-averaging makes the representation approximately invariant to when effects occur within the recording
	\item \textbf{Duration variations}: Since TRR aggregates over entire signal, representations for short and long samples of the same effect are similar
\end{itemize}

These invariances allow TRR to focus on \textit{what} (effect types and their interactions) rather than \textit{when} and \textit{how loud}.

\subsubsection{Uncertainty-Aware Fusion Benefits}

The uncertainty-aware fusion mechanism provides several advantages:

\paragraph{Automatic Modality Selection}. The system automatically adapts to the reliability of each modality without manual tuning:
\begin{itemize}
	\item For \textbf{specific, well-formed prompts} (e.g., ``warm vintage breakup with slight compression''), text retrieval produces a peaked similarity distribution with low entropy, resulting in high text weight
	\item For \textbf{vague or ambiguous prompts} (e.g., ``make it sound good''), text retrieval produces a flat similarity distribution with high entropy, resulting in low text weight
	\item For \textbf{high-quality audio references}, audio retrieval produces confident matches with low uncertainty, receiving high weight
\end{itemize}

\paragraph{Graceful Degradation}. When one modality is unavailable or unreliable, the system gracefully falls back:
\begin{itemize}
	\item \textbf{Missing audio}: The system operates in text-only mode with $w_{\text{audio}} = 0$
	\item \textbf{Poor audio quality}: The system automatically down-weights the unreliable audio modality
	\item \textbf{Missing text}: The system operates in audio-only mode with $w_{\text{text}} = 0$
\end{itemize}

\subsubsection{Constraint Satisfaction Mechanisms}

The constraint validation and repair module ensures parameter usability through:

\paragraph{Immediate Feasibility Checking}. Every generated parameter vector is validated before being sent to the DSP engine:
\begin{itemize}
	\item Range constraints prevent plugin specification violations
	\item Dependency constraints prevent logical inconsistencies (e.g., attack > release)
	\item Activation constraints ensure inactive effects have sensible defaults
\end{itemize}

\paragraph{Retrieval-Guided Repair}. When constraints are violated, the repair mechanism uses retrieved presets as guidance:
\begin{itemize}
	\item Range violations are fixed via simple clipping to valid bounds
	\item Dependency violations are resolved by interpolating with the nearest valid retrieved preset
	\item This preserves as much of the LLM's generation as possible while ensuring feasibility
\end{itemize}

\subsection{Broader Applicability}

The principles behind Audio-Agent extend beyond guitar effects to other domains:

\subsubsection{Other Audio Processing Domains}

\paragraph{Mixing and Mastering}. Similar retrieval-grounded approaches could apply to:
\begin{itemize}
	\item \textbf{Mix balance}: Retrieving reference mixes with similar density and clarity profiles
	\item \textbf{EQ matching}: Using texture features to match equalization curves
	\item \textbf{Dynamic processing}: Retrieving compressor settings based on transient characteristics
\end{itemize}

\paragraph{Synthesis Parameters}. For synthesis instruments (e.g., synthesizers, drum machines), retrieval could help:
\begin{itemize}
	\item \textbf{Preset discovery}: Finding sounds that match text descriptions
	\item \textbf{Sound design}: Adapting retrieved presets for specific contexts
	\item \textbf{Style transfer}: Transferring timbral characteristics across instruments
\end{itemize}

\subsubsection{Non-Audio Domains}

The dual-modal retrieval with uncertainty-aware fusion paradigm generalizes to:
\begin{itemize}
	\item \textbf{Video effects}: Text + video frame retrieval for color grading, motion effects
	\item \textbf{Image processing}: Text + image texture retrieval for filter selection
	\item \textbf{Haptic feedback}: Text + haptic pattern retrieval for tactile interaction design
\end{itemize}

In each case, the core idea remains: use retrieval to provide concrete examples and feasibility constraints, use fusion to handle uncertainty across modalities, and use structured reasoning to generate outputs that respect domain-specific constraints.

\subsection{Detailed Limitations Analysis}

\paragraph{Database Coverage}. The effectiveness of retrieval-based approaches fundamentally depends on the diversity and quality of the preset database. Limitations include:
\begin{itemize}
	\item \textbf{Missing effect combinations}: Out-of-distribution combinations cannot be retrieved
	\item \textbf{Genre gaps}: Under-represented genres (e.g., experimental electronic) have poor retrieval quality
	\item \textbf{Temporal drift}: Database may not reflect current trends in production techniques
\end{itemize}
\textit{Mitigation}: Continuous expansion of database through crowd-sourcing and active learning.

\paragraph{Parameter Identifiability}. Multiple parameter configurations can produce perceptually similar sounds, creating an inherent ambiguity:
\begin{itemize}
	\item \textbf{Many-to-one mapping}: Different parameter vectors yield similar timbres
	\item \textbf{Subjectivity}: What constitutes a ``good'' sound varies by context and preference
	\item \textbf{Interaction effects}: Parameter combinations can have non-intuitive perceptual results
\end{itemize}
\textit{Mitigation}: Incorporate user feedback and preference learning to personalize recommendations.

\paragraph{Computational Constraints}. Real-time deployment imposes practical limitations:
\begin{itemize}
	\item \textbf{Latency requirements}: Audio processing must complete within tight time budgets (< 10ms)
	\item \textbf{Memory constraints}: Embedding caches and retrieval structures must fit in limited memory
	\item \textbf{Scalability}: Larger databases increase retrieval time linearly without careful indexing
\end{itemize}
\textit{Mitigation}: Use efficient indexing (e.g., FAISS), aggressive caching, and approximate nearest neighbors.

\paragraph{Evaluation Limitations}. Our current evaluation has several weaknesses:
\begin{itemize}
	\item \textbf{Objective metrics only}: We use parameter distance as a proxy for perceptual quality, which may not correlate perfectly
	\item \textbf{Limited dataset}: 50 guitar tones may not capture full diversity of use cases
	\item \textbf{No listening tests}: Controlled listening studies are needed to validate perceptual quality
\end{itemize}
\textit{Mitigation}: Conduct MUSHRA-style listening tests with expert and non-expert participants.
